% Options for packages loaded elsewhere
\PassOptionsToPackage{unicode}{hyperref}
\PassOptionsToPackage{hyphens}{url}
\PassOptionsToPackage{dvipsnames,svgnames*,x11names*}{xcolor}
%
\documentclass[
  16pt,
]{krantz}
\usepackage{amsmath,amssymb}
\usepackage{lmodern}
\usepackage{ifxetex,ifluatex}
\ifnum 0\ifxetex 1\fi\ifluatex 1\fi=0 % if pdftex
  \usepackage[T1]{fontenc}
  \usepackage[utf8]{inputenc}
  \usepackage{textcomp} % provide euro and other symbols
\else % if luatex or xetex
  \usepackage{unicode-math}
  \defaultfontfeatures{Scale=MatchLowercase}
  \defaultfontfeatures[\rmfamily]{Ligatures=TeX,Scale=1}
  \setmonofont[Scale=0.7]{Source Code Pro}
\fi
% Use upquote if available, for straight quotes in verbatim environments
\IfFileExists{upquote.sty}{\usepackage{upquote}}{}
\IfFileExists{microtype.sty}{% use microtype if available
  \usepackage[]{microtype}
  \UseMicrotypeSet[protrusion]{basicmath} % disable protrusion for tt fonts
}{}
\makeatletter
\@ifundefined{KOMAClassName}{% if non-KOMA class
  \IfFileExists{parskip.sty}{%
    \usepackage{parskip}
  }{% else
    \setlength{\parindent}{0pt}
    \setlength{\parskip}{6pt plus 2pt minus 1pt}}
}{% if KOMA class
  \KOMAoptions{parskip=half}}
\makeatother
\usepackage{xcolor}
\IfFileExists{xurl.sty}{\usepackage{xurl}}{} % add URL line breaks if available
\IfFileExists{bookmark.sty}{\usepackage{bookmark}}{\usepackage{hyperref}}
\hypersetup{
  pdftitle={Matemáticas básicas},
  pdfauthor={Ricardo Michel MALLQUI BAÑOS},
  colorlinks=true,
  linkcolor=Maroon,
  filecolor=Maroon,
  citecolor=Blue,
  urlcolor=Blue,
  pdfcreator={LaTeX via pandoc}}
\urlstyle{same} % disable monospaced font for URLs
\usepackage{longtable,booktabs,array}
\usepackage{calc} % for calculating minipage widths
% Correct order of tables after \paragraph or \subparagraph
\usepackage{etoolbox}
\makeatletter
\patchcmd\longtable{\par}{\if@noskipsec\mbox{}\fi\par}{}{}
\makeatother
% Allow footnotes in longtable head/foot
\IfFileExists{footnotehyper.sty}{\usepackage{footnotehyper}}{\usepackage{footnote}}
\makesavenoteenv{longtable}
\setlength{\emergencystretch}{3em} % prevent overfull lines
\providecommand{\tightlist}{%
  \setlength{\itemsep}{0pt}\setlength{\parskip}{0pt}}
\setcounter{secnumdepth}{5}
\usepackage[spanish,es-lcroman]{babel}
\usepackage{booktabs}
\usepackage{mathtools}
\usepackage{graphicx}
\usepackage{amsmath}
\usepackage{makeidx}
\makeindex
%\usepackage{showframe}
%\usepackage[a4paper]{geometry}
%\geometry{verbose,tmargin=3cm,bmargin=3cm,lmargin=3.5cm,rmargin=3cm}
\setlength\parindent{23pt}
\usepackage[singlelinecheck=off]{caption}

\usepackage{times}
\renewcommand{\rmdefault}{ptm}
%\usepackage[lite,subscriptcorrection,nofontinfo,zswash]{mtpro2}

\usepackage{graphicx}

% Determine if the image is too wide for the page.
\makeatletter
\def\ScaleIfNeeded{%
  \ifdim\Gin@nat@width>\linewidth
    \linewidth
  \else
    \Gin@nat@width
  \fi
}
\makeatother

% Resize figures that are too wide for the page.
%\let\oldincludegraphics\includegraphics
%\renewcommand\includegraphics[2][]{%
%  \oldincludegraphics[scale=0.85]{#2}
%}

\usepackage{amsthm}
\makeatletter
\def\thm@space@setup{%
  \thm@preskip=8pt plus 2pt minus 4pt
  \thm@postskip=\thm@preskip
}
\makeatother

\flushbottom 

\frontmatter

\ifluatex
  \usepackage{selnolig}  % disable illegal ligatures
\fi
\usepackage[]{natbib}
\bibliographystyle{apalike}

\title{Matemáticas básicas}
\author{Ricardo Michel MALLQUI BAÑOS}
\date{2021-06-16}

\usepackage{amsthm}
\newtheorem{theorem}{Teorema}[chapter]
\newtheorem{lemma}{Lema}[chapter]
\newtheorem{corollary}{Corolario}[chapter]
\newtheorem{proposition}{Proposición}[chapter]
\newtheorem{conjecture}{Conjectura}[chapter]
\theoremstyle{definition}
\newtheorem{definition}{Definición}[chapter]
\theoremstyle{definition}
\newtheorem{example}{Ejemplo}[chapter]
\theoremstyle{definition}
\newtheorem{exercise}{Ejercicio}[chapter]
\theoremstyle{definition}
\newtheorem{hypothesis}{Hypothesis}[chapter]
\theoremstyle{remark}
\newtheorem*{remark}{Observación}
\newtheorem*{solution}{Solución}
\begin{document}
\maketitle

%\cleardoublepage\newpage\thispagestyle{empty}\null
%\cleardoublepage\newpage\thispagestyle{empty}\null
%\cleardoublepage\newpage
\thispagestyle{empty}
\begin{flushright}
Universidad Nacional San Cristobal de Huamanga

Fisart.cf

Agradecimento a los estudiantes de la ESFAPA FGPA

A la UNSCH

\includegraphics[height=3cm]{U.pdf}
\end{flushright}
%\setlength{\abovedisplayskip}{-5pt}
%\setlength{\abovedisplayshortskip}{-5pt}

{
\hypersetup{linkcolor=}
\setcounter{tocdepth}{2}
\tableofcontents
}
\listoftables
\listoffigures
\newcommand{\N}{\mathbb{N}}
\newcommand{\R}{\mathbb{R}}
\newcommand{\CC}{\mathbb{C}}
\newcommand{\I}{\mathbb{I}}
\newcommand{\f}{\mathbb{f}}
\newcommand{\X}{\mathbb{X}}
\newcommand{\D}{\mathbb{D}}
\newcommand{\Z}{\mathbb{Z}}
\newcommand{\Q}{\mathbb{Q}}
\newcommand{\norm}[1]{\left\Vert#1\right\Vert}
\newcommand{\abs}[1]{\left\vert#1\right\vert}
\newcommand{\set}[1]{\left\{#1\right\}}
\newcommand{\seq}[1]{\left<#1\right>}
\newcommand{\co}[1]{\left[#1\right]}
\newcommand{\cc}[1]{\left(#1\right)}
\newcommand{\J}{\mathcal{J}}
\newcommand{\K}{\mathcal{K}}
\newcommand{\M}{\mathcal{M}}
\newcommand{\F}{\mathcal{F}}

\hypertarget{resumen}{%
\chapter*{Resumen}\label{resumen}}


www.

\hypertarget{introducciuxf3n}{%
\chapter*{Introducción}\label{introducciuxf3n}}


www.

\mainmatter

\hypertarget{logica}{%
\chapter{Logica}\label{logica}}

www.

\hypertarget{conjuntos}{%
\chapter{Conjuntos}\label{conjuntos}}

\begin{definition}[Conjunto]
\protect\hypertarget{def:conjunto}{}{\label{def:conjunto} \iffalse (Conjunto) \fi{} }Es una coleccion de elementos con caractersiticas similares
\end{definition}

\begin{definition}[Determinacion de conjuntos]
\protect\hypertarget{def:conjunto2}{}{\label{def:conjunto2} \iffalse (Determinacion de conjuntos) \fi{} }Por extensión y comprensión
\end{definition}

Cornutos universal, vacio, unitario

\[
\mathbb{N}, \mathbb{Z}, \mathbb{Q}, \mathbb{I}, \mathbb{R}, \mathbb{C}
\]

\[
\phi=\left\{x/x\neq x\right\}
\]

\[
A=\left\{a\right\}
\]

\hypertarget{funciuxf3n-proposicional-y-cuantificadores}{%
\section{Función proposicional y cuantificadores}\label{funciuxf3n-proposicional-y-cuantificadores}}

\hypertarget{funciuxf3n-proposicional}{%
\subsection{Función proposicional}\label{funciuxf3n-proposicional}}

\begin{definition}[Función proposicional]
\protect\hypertarget{def:proposicional}{}{\label{def:proposicional} \iffalse (Función proposicional) \fi{} }
\end{definition}

Sea \(x\) una variable \(P(x)\) un enunciado, \(P(x)\) es una \textbf{\emph{función proposicional}} si al sustituir la variable con una constante este se convierte en una proposición.

Por ejemplo \(P(x)\): \(x\) es un numero par

Al conjunto de todos lo valores de \(x\) se denomina \emph{domino de la variable}

\hypertarget{cuantificadores}{%
\subsection{Cuantificadores}\label{cuantificadores}}

\begin{definition}[Cuantificador existencial]
\protect\hypertarget{def:existencial}{}{\label{def:existencial} \iffalse (Cuantificador existencial) \fi{} }
\end{definition}

\[
\exists
\]

Es una generalización de la disyunción Inclusiva. Por ello, es verdadero cuando al menos un valor de \(x\) perteneciente al Dominio de \(A\), es Verdadero. Se denota; \(\exists x / P (x)\) Se lee: ``Existe al menos un \(x\)'', ``Algunos \(x\)''," Hay \(x\)``,''Existe un \(x\)", etc.

\begin{definition}[Cuantificador universal]
\protect\hypertarget{def:universal}{}{\label{def:universal} \iffalse (Cuantificador universal) \fi{} }
\end{definition}

\[
\forall
\]

Es una generalización de la \emph{conjunción}. Debido a esto es verdadero cuando todos los valores de \(x\) que pertenecen al Dominio de \(A\) son Verdaderos. Se denota: \(\forall x ; p(x)\) Se lee: ``Para Todo \(x\)'', ``Para cada \(x\)'', ``Todos (as) las \(x\)'', ``Todo \(x\)''.

Sea \(A=\left\{1,2,3,4,5\right\}\) y la función proposicional \(3x-2<12\) entonces las proposiciones

\begin{enumerate}
\def\labelenumi{\arabic{enumi}.}
\tightlist
\item
  \(\forall x\in A;3x-1<14\)
\item
  \(\exists\; x\in A|3x-2<12\)
\end{enumerate}

son falsa y verdadera y se escribe

\begin{definition}[Proposición universal]
\protect\hypertarget{def:universal2}{}{\label{def:universal2} \iffalse (Proposición universal) \fi{} }
\end{definition}

Una proposición universal es aquella que está provista de un cuantificador universal, y tiene la forma: \[\forall x\in A:p(x)\]

\begin{definition}[Proposición existencial]
\protect\hypertarget{def:existencial2}{}{\label{def:existencial2} \iffalse (Proposición existencial) \fi{} }
\end{definition}

Una proposición universal es aquella que está provista de un cuantificador existencial, y tiene la forma: \[\exists x\in A:p(x)\]

\hypertarget{negaciuxf3n-de-las-proposiciones-universal-y-existencial}{%
\subsection{Negación de las proposiciones universal y existencial}\label{negaciuxf3n-de-las-proposiciones-universal-y-existencial}}

Cambiando el cuantificador universal por el cuantificador existencial, o viceversa

es decir \[\sim[\exists x; P(x)]\equiv\forall x;\sim P(x)\]

\[
\sim\left[\forall  x; P(x)\right]\equiv\exists\; x;\sim P(x)
\]

La negación del \textbf{\emph{cuantificador universal}} es equivalente a la \emph{afirmación de un cuantificador existencial} respecto de la \textbf{\emph{función proposicional negada}.}

La negación de un \textbf{\emph{cuantificador existencial}} es equivalente a la \emph{afirmación de un cuantificador universal} respecto de la \textbf{\emph{función proposicional negada}}.

\begin{example}
\protect\hypertarget{exm:unnamed-chunk-1}{}{\label{exm:unnamed-chunk-1} }
\end{example}

Dada la proposición: ``Si todos los números primos son impa­res, los números positivos son mayores que -1''

\begin{itemize}
\tightlist
\item
  Expresarla simbólicamente
\item
  Negar oracianalmente la proposición
\end{itemize}

\begin{solution}
\iffalse{} {Solución. } \fi{}
\end{solution}

Sea \(p(x):\) números primos impares y \(q(x):\) numeros positivos mayores que -1

\begin{itemize}
\item
  \(\forall x:[p(x)\rightarrow q(x)]\)
\item
\item
\end{itemize}

\begin{example}
\protect\hypertarget{exm:unnamed-chunk-3}{}{\label{exm:unnamed-chunk-3} }
\end{example}

Dado el conjunto \(A=\left\{x\in\mathbb{N}:-14<x<27\right\}\). Hallar el valor de verdad de \[
s=[(\sim p\wedge \sim q)\rightarrow(\sim q\wedge \sim r)]\leftrightarrow(\sim p\vee r) 
\] si \(p=(\forall x\in A, \exists y\in A, \forall z\in A)[x^2-z^2>y^2]\), \(q=(\exists y\in A, \forall z\in A, \exists x \in A)[2x-4y<-z]\) y \(q=(\forall z\in A, \exists x\in A, \forall y \in A)[3x^2-z^2>y]\)

\begin{solution}
\iffalse{} {Solución. } \fi{}
\end{solution}

\(A\left\{1,2,3,\ldots,26\right\}\) luego el valor de \(\text{V}(p)=F\), \(\text{V}(q)=F\) y \(\text{V}(r)=V\) por lo tanto

\[
\begin{aligned}
\text{V}(s)&=[(\sim p\wedge \sim q)\Longrightarrow(\sim q\wedge \sim r)]\Longleftrightarrow(\sim p\vee r)\\
&=[(V\wedge F)\Longrightarrow(F\wedge F)]\Longleftrightarrow(V\vee V)\\
&=[F\Longrightarrow F]\Longleftrightarrow V\\
&=V
\end{aligned}
\]

\hypertarget{operaciones-entre-conjuntos}{%
\section{Operaciones entre conjuntos}\label{operaciones-entre-conjuntos}}

\hypertarget{relaciones-entre-conjuntos-conjuntos-iguales.-conjuntos-equivalentes}{%
\section{Relaciones entre Conjuntos: Conjuntos Iguales. Conjuntos equivalentes}\label{relaciones-entre-conjuntos-conjuntos-iguales.-conjuntos-equivalentes}}

\hypertarget{representaciuxf3n-gruxe1fica-de-los-conjuntos}{%
\section{Representación Gráfica de los Conjuntos}\label{representaciuxf3n-gruxe1fica-de-los-conjuntos}}

Diagrama de euler

\hypertarget{uniuxf3n-de-conjuntos}{%
\section{Unión de Conjuntos}\label{uniuxf3n-de-conjuntos}}

\hypertarget{intersecciuxf3n-de-conjuntos.-propiedades}{%
\section{Intersección de Conjuntos. Propiedades}\label{intersecciuxf3n-de-conjuntos.-propiedades}}

\hypertarget{distributivas-de-la-uniuxf3n-e-intersecciuxf3n}{%
\section{Distributivas de la Unión e Intersección}\label{distributivas-de-la-uniuxf3n-e-intersecciuxf3n}}

\hypertarget{leyes-de-absorciuxf3n}{%
\section{Leyes de Absorción}\label{leyes-de-absorciuxf3n}}

\hypertarget{diferencia-de-conjuntos.}{%
\section{Diferencia de Conjuntos.}\label{diferencia-de-conjuntos.}}

\hypertarget{complemento-de-un-conjunto.-propiedades}{%
\section{Complemento de un Conjunto. Propiedades}\label{complemento-de-un-conjunto.-propiedades}}

\hypertarget{diferencia-simuxe9trica.}{%
\section{Diferencia Simétrica.}\label{diferencia-simuxe9trica.}}

\hypertarget{nuxfamero-de-elementos-de-un-conjunto.-propiedades}{%
\section{Número de elementos de un Conjunto. Propiedades}\label{nuxfamero-de-elementos-de-un-conjunto.-propiedades}}

\hypertarget{funciones-y-relaciones}{%
\chapter{Funciones y relaciones}\label{funciones-y-relaciones}}

\hypertarget{numeros-reales}{%
\chapter{Numeros reales}\label{numeros-reales}}

\hypertarget{funciones-exponenciales-logaruxedtmicas}{%
\chapter{Funciones exponenciales logarítmicas}\label{funciones-exponenciales-logaruxedtmicas}}

\hypertarget{inducciuxf3n-matemuxe1tica}{%
\chapter{Inducción matemática}\label{inducciuxf3n-matemuxe1tica}}

\hypertarget{suceciones}{%
\chapter{Suceciones}\label{suceciones}}

\hypertarget{nuxfameros-complejos}{%
\chapter{Números complejos}\label{nuxfameros-complejos}}

\hypertarget{polinomios}{%
\chapter{Polinomios}\label{polinomios}}

\hypertarget{appendix-apendice}{%
\appendix \addcontentsline{toc}{chapter}{\appendixname}}


Temas de reforzamiento o conocimientos preliminares que son necesarias para entender el contenido.

\hypertarget{trasformaciones}{%
\chapter{Trasformaciones}\label{trasformaciones}}

  \bibliography{book.bib,packages.bib}

\printindex

\end{document}
